\documentclass[twoside,english,aps,showpacs,pra,superscriptaddress]{revtex4-2}
\usepackage{amsmath}
\usepackage{amssymb}
\usepackage{amsfonts}
\usepackage{amsthm}
\usepackage{setspace}
\usepackage{float}
\usepackage{subfigure}
\usepackage{graphicx}
\usepackage{subfigure}
\usepackage{braket}
\usepackage{url}
\usepackage{cancel}
\usepackage{extarrows}
\usepackage{color}
\usepackage[plainpages = false, pdfpagelabels, 
                % pdfpagelayout = useoutlines,
                 bookmarks,
                 bookmarksopen = true,
                 bookmarksnumbered = true,
                 breaklinks = true,
                 linktocpage,
                 %pagebackref,
                 colorlinks = true,
                 linkcolor = blue,
                 urlcolor  = blue,
                 citecolor = blue,
                 anchorcolor = green,
                 hyperindex = true,
                 hyperfigures
                 ]{hyperref} 
\renewcommand{\d}{\mathrm{d}}
\newcommand{\D}{\mathrm{D}}
\newcommand{\e}{\mathrm{e}}
\newcommand{\ii}{\mathrm{i}}
\DeclareMathOperator{\Tr}{Tr}
\DeclareMathOperator{\rank}{rank}
\DeclareMathOperator{\tr}{tr}
\DeclareMathOperator{\im}{im}
\DeclareMathOperator{\sech}{sech}
\DeclareMathOperator{\wt}{wt}
\DeclareMathOperator*{\argmin}{argmin}
\renewcommand{\vec}{\boldsymbol}
\renewcommand{\abstractname}{}

\newcommand{\sket}[1]{\left.\left| {#1} \right\rangle\right\rangle}
\newcommand{\sbra}[1]{\left\langle\left\langle {#1} \right|\right.}
\newcommand{\sbraket}[1]{\left\langle\left\langle {#1}\right\rangle\right\rangle}
\newcommand{\norm}[1]{\left| \left| {#1}\right|\right|}

\newtheorem{lemma}{Lemma}
\newtheorem{theorem}{Theorem}
\newtheorem{proposition}{Proposition}
\newtheorem{definition}{Definition}

\setcounter{equation}{0}
\setcounter{figure}{0}
\setcounter{table}{0}
\setcounter{page}{1}
\makeatletter


\begin{document}
\title{Single-shot State Preparation}  
\maketitle
\section{motivation}
Confinement is a weaker condition compared with soundness to guarantee single-shot decoding \cite{quintavalleSingleShotErrorCorrection2021}. It could be satisfied by code with no check redundancy, e.g. quantum expander code \cite{fawziConstantOverheadQuantum2018b}. Such codes exhibit no thermodynamic stability but still support single-shot decoding as well as self-correction \cite{hongQuantumMemoryNonzero2025a}. The question is which condition suffices to guarantee measurement preparations of code states in the presence of stochastic readout errors? Does single-shot decoding imply single-shot preparation? Is confinement enough, or we further need soundness/thermodynamic stability?

\section{outline}
Measurement preparations of code states effectively induce a thermal state in the presence of stochastic readout errors.
Consider a CSS code $\mathcal{K}$, which is defined by a length-3 chain complex,
\begin{equation}
\begin{aligned}
    &C_2 \xrightarrow{\partial_2 = H_Z^T} C_1 \xrightarrow{\partial_1 = H_X}C_0,\\
    &C^2 \xleftarrow{\delta_1 = H_Z} C^1 \xleftarrow{\delta_0 = H_X^T}C^0.
\end{aligned}
\end{equation}
Here $H_X$ and $H_Z$ are party check matrices. The logical $Z$ operators are given by $H_1 = Z_1/B_1$, where $Z_1=\ker \partial_1$, $B_1=\im \partial_2$, and logical $X$ operators are given by $H^1 = Z^1/B^1$, where $Z^1=\ker \delta_0$, $B_1=\im \delta_1$. The basis of $C_2$ is denoted as $e_{2,i}$, and each of them induces a $Z$ stabilizer generator $g_{Z,i} = Z_{H_Z^T e_{2,i}} = \prod_{e_1 \in H_Z^T e_{2,i}} Z_{e_1}$. $X$ stabilizer generator $g_X$ can be defined similarly. 

Suppose that we start from the product $+$ state $\ket{+}^{\otimes n}$ and measure all $Z$ checks (basis of $C_2$) to prepare logical $+$ state $[h^1] = 0$. The measurement of each generator $g_Z$ is given by the following channel
\begin{equation}
    ((1-q) \ket{0}_c\bra{0}_c + q\ket{1}_c \bra{1}_c)\Pi_+ \rho_0 \Pi_+ + ((1-q) \ket{1}_c\bra{1}_c + q\ket{0}_c \bra{0}_c)\Pi_- \rho_0 \Pi_-, \quad \Pi_{\pm} = \frac{1\pm g_Z}{2},
\end{equation}
where the measurement records are stored in the classical register.
Since the initial product state can be expanded in the stabilizer basis as 
\begin{equation}
    \rho_0 = \ket{+}^{\otimes n} \bra{+}^{\otimes n} = \bigotimes_{g_Z'} \frac{1}{2}\left(\ket{g_Z'=+}  + \ket{g_Z'=-}\right)\left(\bra{g_Z'=+}  + \bra{g_Z'=-}\right) \bigotimes_{g_X'} \ket{g_X'=+} \bra{g_X'=+} \bigotimes_{L_X} \ket{L_X=+}\bra{L_X=+},
\end{equation}
where ${g_{Z}'}$ (${g_{X}'}$) is the full-rank generator set associated with $C_2/\ker H_Z^T$ ($C^0/\ker H_X^T$) without redundancy, and $L_Z$ ($L_X$) denotes the logical $Z$ ($X$) generator associated with basis of $H_1$ ($H^1$). Since stabilizer measurements commute with each other, we can apply all $g_Z'$ measurements to $\rho_0$ first, such that each $g_Z'$ sector becomes
\begin{equation}
    \begin{aligned}
        &  \frac{1}{2}\sum_{s_{g_Z'} = \pm }((1-q) \ket{s_{g_Z'}}_c\bra{s_{g_Z'}}_c + q\ket{-s_{g_Z'}}_c \bra{-s_{g_Z'}}_c) \otimes \Pi_{s_{g_Z'}} \left(\ket{g_Z'=+}  + \ket{g_Z'=-}\right)\left(\bra{g_Z'=+}  + \bra{g_Z'=-}\right)\Pi_{s_{g_Z'}}\\
        & =  \frac{1}{2}\sum_{s_{g_Z'} = \pm }((1-q) \ket{s_{g_Z'}}_c\bra{s_{g_Z'}}_c + q\ket{-s_{g_Z'}}_c \bra{-s_{g_Z'}}_c)\otimes\Pi_{s_{g_Z'}} \\
        & =  \frac{1}{2}\sum_{s_{g_Z'} = \pm } \ket{s_{g_Z'}}_c\bra{s_{g_Z'}}_c \otimes \frac{1+(1-2q) s_{g_Z'} g_Z'}{2} \propto \sum_{s_{g_Z'} = \pm } \ket{s_{g_Z'}}_c\bra{s_{g_Z'}}_c \otimes \exp\left(\beta s_{g_Z'} g_Z'\right),
    \end{aligned}
\end{equation}
where 
\begin{equation}
    \tanh \beta = 1-2q.
\end{equation}
Then we apply the remaining redundant $g_Z$ measurements. Since they commute with $ \exp\left(\beta s_{g_Z'} g_Z'\right)$ because every $g_Z$ is a product of several $g_Z'$, we have the post measurement state
\begin{equation}
    \rho_Q = \frac{1}{N}\sum_{\{s_{g_Z} \} } \ket{\{s_{g_Z}\}}_c\bra{\{s_{g_Z}\}}_c \otimes \exp\left(\beta \sum_{g_Z} s_{g_Z} g_Z\right) \bigotimes_{g_X'} \ket{g_X'=+} \bra{g_X'=+} \bigotimes_{L_X} \ket{L_X=+}\bra{L_X=+}.
\end{equation}
Denote $r_{X/Z}' = \rank H_{X/Z}$, $r_{X} = \dim C^0$, $r_Z = \dim C_2$. We have $n= r_Z'+r_X'+k$.
The normalization constant is given by
\begin{equation}
\begin{aligned}
    &N = \sum_{\{s_{g_Z} \} } \tr_Q \left(  \exp\left(\beta \sum_{g_Z} s_{g_Z} g_Z\right) \bigotimes_{g_X'} \ket{g_X'=+} \bra{g_X'=+} \bigotimes_{L_X} \ket{L_X=+}\bra{L_X=+}\right) \\
    & = \frac{1}{2^{r_X'+k}}\sum_{\{s_{g_Z} \} } \tr_Q \exp\left(\beta \sum_{g_Z} s_{g_Z} g_Z\right) \\
    & = \frac{2^n}{2^{r_X'+k}} (2\cosh\beta)^{r_Z} = 2^{r_Z'} (2\cosh\beta)^{r_Z}.
\end{aligned}
\end{equation}
The probability of each measurement outcome configuration is given by 
\begin{equation}
\begin{aligned}
    & \Pr(\{s_{g_Z}\}) =  \frac{1}{N} \tr_Q \left(  \exp\left(\beta \sum_{g_Z} s_{g_Z} g_Z\right) \bigotimes_{g_X'} \ket{g_X'=+} \bra{g_X'=+} \bigotimes_{L_X} \ket{L_X=+}\bra{L_X=+}\right)\\
    & = \frac{1}{2^{r_Z'} (2\cosh\beta)^{r_Z}} \mathcal{Z}(\{s_{g_Z}\}), 
\end{aligned}
\end{equation}
where
\begin{equation}
    \mathcal{Z}(\{s_{g_Z}\}) = \frac{1}{2^{r_X'+k}}\tr_Q \exp\left(\beta \sum_{g_Z} s_{g_Z} g_Z\right) 
\end{equation}
is the partition function. So we have
\begin{equation}
    \rho_Q = \sum_{\{s_{g_Z} \} } \Pr(\{s_{g_Z}\}) \ket{\{s_{g_Z}\}}_c\bra{\{s_{g_Z}\}}_c \bigotimes \frac{1}{\mathcal{Z}(\{s_{g_Z}\})}\exp\left(\beta \sum_{g_Z} s_{g_Z} g_Z\right) \bigotimes_{g_X'} \ket{g_X'=+} \bra{g_X'=+} \bigotimes_{L_X} \ket{L_X=+}\bra{L_X=+}.
\end{equation}
The common state preparation procedure involves a round of decoding to flip the syndrome towards the ideal logical state $g_Z=+$. The decoder applies certain correction channel $\mathcal{K}_{\{s_{g_Z}\}}$ to the system with respect to the measurement record $\{s_{g_Z}\}$. Normally the correction only involves Pauli $X$ operators, so it only applies to the $g_Z$ sector, such that  
\begin{equation}
    \mathcal{K}(\rho_Q) = \sum_{\{s_{g_Z}\} } \Pr(\{s_{g_Z}\} ) \ket{\{s_{g_Z}\}}_c\bra{\{s_{g_Z}\}}_c \bigotimes \mathcal{K}_{\{s_{g_Z}\}} \left(\frac{\exp\left(\beta \sum_{g_Z} s_{g_Z} g_Z\right) }{\mathcal{Z}(\{s_{g_Z}\})} \right)\bigotimes_{g_X'} \ket{g_X'=+} \bra{g_X'=+} \bigotimes_{L_X} \ket{L_X=+}\bra{L_X=+}.
\end{equation}
The next step is encode logical information in the above state. Physically it is done by logical operators but mathematically this could be described by the encoding channel
\begin{equation}
    \mathcal{E}'(\rho_L) = \sum_{\{s_{g_Z}\} } \Pr(\{s_{g_Z}\} ) \ket{\{s_{g_Z}\}}_c\bra{\{s_{g_Z}\}}_c \bigotimes \mathcal{K}_{\{s_{g_Z}\}} \left(\frac{\exp\left(\beta \sum_{g_Z} s_{g_Z} g_Z\right) }{\mathcal{Z}(\{s_{g_Z}\})} \right)\bigotimes_{g_X'} \ket{g_X'=+} \bra{g_X'=+} \bigotimes \rho_L.
\end{equation}
In comparison, the undecoded encoding channel is
\begin{equation}
    \mathcal{E}(\rho_L) = \sum_{\{s_{g_Z}\} } \Pr(\{s_{g_Z}\} ) \ket{\{s_{g_Z}\}}_c\bra{\{s_{g_Z}\}}_c \bigotimes  \frac{\exp\left(\beta \sum_{g_Z} s_{g_Z} g_Z\right) }{\mathcal{Z}(\{s_{g_Z}\})} \bigotimes_{g_X'} \ket{g_X'=+} \bra{g_X'=+} \bigotimes \rho_L.
\end{equation}
We want to show that the encoding $\mathcal{E}'$ cannot protect quantum information against stochastic $X$ noise,
\begin{equation}
    \mathcal{N} (\rho) = \sum_{c \in C_1} \Pr(c) X_c \rho X_c.
\end{equation}

To this end, we utilize coherent information of the composed channel $I_c(\rho_L,\mathcal{N\circ E'})$. Under the assumption that there are only $O(1)$ check redundancies, $\dim \ker H_Z^T = O(1)$, we need to address the following questions,
\begin{enumerate}
    \item Does $I_c(\rho_L,\mathcal{N\circ E'}) \leq I_c(\rho_L,\mathcal{N\circ E})$ hold?
    \item Does $\lim_{n\rightarrow\infty} I_c(\rho_L,\mathcal{N\circ E}) < k\log 2$ (or $\lim_{n\rightarrow\infty} I_c(\rho_L,\mathcal{N\circ E})/n < \kappa\log 2$ and $\kappa$ is the code rate) hold?
    \item Whether $\lim_{n\rightarrow\infty}  I_c(\rho_L,\mathcal{N\circ E'}) < k \log 2$ (or $\lim_{n\rightarrow\infty} I_c(\rho_L,\mathcal{N\circ E'})/n < \kappa\log 2$) implies $\lim_{n\rightarrow\infty}F_e(\mathcal{R\circ N\circ E'}) < 1$, $\forall\mathcal{R}$, especially when $\mathcal{E}'$ is a mixed state encoding?
\end{enumerate}
If the above three propositions holds, we conclude there is not single-shot preparation threshold without check redundancy.
Additionally, we should address
\begin{enumerate}
    \item Given that 
    \begin{equation}
    \begin{aligned}
        &\rho_Q = \sum_s \Pr(s) \ket{s}_c\bra{s}_c \otimes \frac{(2\cosh\beta)^{r_Z}}{\mathcal{Z}(s)}\sum_{s_e|s_e+s\in \im H_Z}  (1-q)^{r_Z - s_e} q^{s_e} \ket{\sigma_{Z} = s+s_e}\bra{\sigma_{Z} = s+s_e}\\& \otimes \ket{\sigma_{X}' = 0}\bra{\sigma_{X}' = 0} \otimes \ket{[h]=0}\bra{[h]=0},
    \end{aligned}
    \end{equation}
    where we denoted the measurement record as $s\in C_2$ and the measurement error as $s_e \in C_2$, will soundness guarantee single-shot preparation? Note that $\sigma_X' \in C^0/\ker{H_X^T}$ has no redundancy, but $\sigma_Z \in C_2$ could lead to redundant projectors.
    \item A Clifford simulation comparing quantum expander code and $3$-d toric code preparation? Procedure: product state $\rightarrow$ projection preparation $\rightarrow$ decode $\rightarrow$ encode logical information $\rightarrow$ stochastic $X$ noise $\rightarrow$ syndrome measurement $\rightarrow$ decode $\rightarrow$ ... $\rightarrow$ calculate logical error rate.
\end{enumerate}

\section{Proof of proposition 1}
\begin{lemma}
Suppose that $\mathcal{K}$ is a decoding channel that applies Pauli $X$ recovery according to the measurement record, such that the decoded encoding channel $\mathcal{E}' = \mathcal{K} \circ \mathcal{E}$. Then we have
\begin{equation}
    I_c(\rho_L, \mathcal{N} \circ \mathcal{E}') \leq I_c(\rho_L, \mathcal{N} \circ \mathcal{E}).
\end{equation}
\end{lemma}
\begin{proof}
The channel $\mathcal{K}$ reads the measurement record $s$ (stored in the classical register from Eq.~(10)) and applies a recovery operation $U_s$. According to the problem setup, the decoder aims to correct $Z$-stabilizer syndrome violations, and thus the recovery operator $U_s$ consists solely of Pauli $X$ operators.

The noise channel $\mathcal{N}$, defined in Eq.~(14), corresponds to stochastic Pauli $X$ noise. Since any Pauli $X$ operator commutes with any other Pauli $X$ operator ($[X_i, X_j] = 0$), the recovery operators $U_s$ commute with the Kraus operators of the noise channel $\mathcal{N}$. Consequently, the correction channel commutes with the noise channel:
\begin{equation}
    [\mathcal{N}, \mathcal{K}] = 0 \implies \mathcal{N} \circ \mathcal{K} = \mathcal{K} \circ \mathcal{N}.
\end{equation}
By substituting this commutation relation into the composed channel expression, we have:
\begin{equation}
    \mathcal{N} \circ \mathcal{E}' = \mathcal{N} \circ (\mathcal{K} \circ \mathcal{E}) = (\mathcal{N} \circ \mathcal{K}) \circ \mathcal{E} = \mathcal{K} \circ (\mathcal{N} \circ \mathcal{E}).
\end{equation}
The coherent information $I_c$ satisfies the Data Processing Inequality \cite{schumacherQuantumDataProcessing1996}, which states that for any state $\rho$ and quantum channels $\Phi, \Lambda$, the inequality $I_c(\rho, \Lambda \circ \Phi) \leq I_c(\rho, \Phi)$ holds. Identifying $\Phi = \mathcal{N} \circ \mathcal{E}$ as the undecoded noisy channel and $\Lambda = \mathcal{K}$ as the post-processing correction, we obtain:
\begin{equation}
\begin{aligned}
    I_c(\rho_L, \mathcal{N} \circ \mathcal{E}') &= I_c(\rho_L, \mathcal{K} \circ (\mathcal{N} \circ \mathcal{E})) \leq I_c(\rho_L, \mathcal{N} \circ \mathcal{E}).
\end{aligned}
\end{equation}
This result implies that if the correction operations commute with the noise, the decoding step cannot recover the quantum information lost to the environment, confirming the proposition.
\end{proof}

\section{Proof of proposition 3}
\begin{lemma}
    if the coherent information density satisfies
    \begin{equation}
        \lim_{n\rightarrow \infty} I_c(\Gamma, \mathcal{N} \circ \mathcal{E}')/k < \log 2,
    \end{equation}
    where $\Gamma = I/2^k$ is the maximally mixed state. Then for any recovery channel $\mathcal{R}$, we have
    \begin{equation}
        \lim_{n\rightarrow \infty} F_e(\mathcal{R\circ N \circ E', I}) < 1,
    \end{equation}
    where $F_e(\mathcal{R\circ N \circ E', I})$ is the average entanglement fidelity between $\mathcal{R\circ N \circ E}$ and the identity channel $\mathcal{I}$.
\end{lemma}
\begin{proof}
    The proof in \cite{zhaoExtractingErrorThresholds2024} still holds.
    
    The proof relies on the Quantum Fano Inequality (also known as the converse bound for quantum capacity). Let the effective channel after recovery be $\Lambda = \mathcal R \circ \mathcal N \circ \mathcal E'$. The relationship between the coherent information loss and the entanglement fidelity $F_e(\Lambda,\mathcal{I})$ is given by:
\begin{equation} \label{eq:fano}
    S(\Gamma) - I_c(\Gamma, \Lambda) \leq H_b(F_e(\Lambda,\mathcal{I})) + (1 - F_e(\Lambda,\mathcal{I})) \log(2^{2k} - 1),
\end{equation}
where $S(\Gamma) = k \log 2$ is the von Neumann entropy of the source, and $H_b(x) = -x\log x - (1-x)\log(1-x)$ is the binary entropy function.

First, we apply the Data Processing Inequality (DPI) for coherent information. Since the recovery operation $\mathcal R$ is a trace-preserving completely positive (TPCP) map acting after the effective channel $\mathcal N \circ \mathcal E'$, it cannot increase the coherent information:
\begin{equation} \label{eq:dpi}
    I_c(\Gamma, \mathcal R \circ (\mathcal N \circ \mathcal E')) \leq I_c(\Gamma, \mathcal N \circ \mathcal E').
\end{equation}

Combining Eq. (\ref{eq:fano}) and Eq. (\ref{eq:dpi}), we obtain a lower bound on the ``error term'' (RHS of Eq. (\ref{eq:fano})):
\begin{equation}
    S(\Gamma) - I_c(\Gamma, \mathcal N \circ \mathcal E') \leq H_b(F_e(\Lambda,\mathcal{I})) + (1 - F_e(\Lambda,\mathcal{I})) \log(2^{2k} - 1).
\end{equation}

By the hypothesis, there exists a strictly positive deficit $\delta > 0$ such that in the asymptotic limit:
\begin{equation}
    \lim_{n\to\infty} \left(  \log 2 - I_c(\Gamma, \mathcal N \circ \mathcal E') /k\right) = \delta>0.
\end{equation}
Consequently, the fidelity $F_e$ must satisfy:
\begin{equation}
    \delta \leq \lim_{n\to\infty} \frac{1}{k}\left(H_b(F_e) + (1 - F_e) \log(2^{2k} - 1)\right).
\end{equation}
Consider the function $f(x) = (H_b(x) + (1-x)\log(2^{2k} - 1))/k$. We observe that $f(x)$ is continuous for $x \in [0, 1]$ and $f(1) = 0$ and is monotonically decreasing for $x\in [1/2,1]$. Since $\delta > 0$, it implies that $F_e$ cannot converge to 1. Notice that this conclusion holds even for $k=\omega(1)$, since in that case $f(x) \rightarrow (1-x) 2\log2$.

\end{proof}


\section{proof of proposition 2}
To prove $\lim_{n\rightarrow\infty} I_c(\rho_L, \mathcal{N} \circ \mathcal{E})/k < \log 2 $, we notice that 
\begin{equation}
    \mathcal{E} = \sum_{s} \Pr(s) \ket{s}_c\bra{s}_c \otimes \mathcal{E}_s,
\end{equation}
where 
\begin{equation}
     \mathcal{E}_s(\rho_L) = \frac{\exp\left(\beta \sum_{g_Z} s_{g_Z} g_Z\right) }{\mathcal{Z}(\{s_{g_Z}\})} \bigotimes_{g_X'} \ket{g_X'=+} \bra{g_X'=+} \bigotimes \rho_L,
\end{equation}
such that
\begin{equation}
    I_c(\rho_L, \mathcal{N} \circ \mathcal{E}) = \sum_{s} \Pr(s) I_c(\rho_L, \mathcal{N} \circ \mathcal{E}_s).
\end{equation}
So it suffice to prove $\lim_{n\rightarrow\infty} I_c(\rho_L, \mathcal{N} \circ \mathcal{E}_s)/k < \log 2 $.
To do so, we note that 
\begin{equation}
     \exp\left(\beta \sum_{g_Z} s_{g_Z} g_Z\right) = \cosh^{r_Z}\beta \sum_{\sigma_Z \in C_2} (\tanh \beta)^{|\sigma_Z|} s_{\sigma_Z} Z_{H_Z^T \sigma_Z},
\end{equation}
where \begin{equation}
    Z_{H_Z^T \sigma_Z} = \prod_{e_1 \in H_Z^T \sigma_Z} Z_{e_1}, \quad s_{\sigma_Z} = \prod _{e_2 \in \sigma_Z} s_{e_2},
\end{equation}
each $Z$ generator $g_Z$ is associated with a basis vector $e_2$ of $C_2$. So
\begin{equation}
    \mathcal{Z}(\{s_{g_Z}\}) = \frac{1}{2^{r_X'+k}}\tr_Q \exp\left(\beta \sum_{g_Z} s_{g_Z} g_Z\right) = 2^{r_Z'} \cosh^{r_Z}\beta \sum_{\sigma_Z \in \ker H_Z^T} (\tanh \beta)^{|\sigma_Z|} s_{\sigma_Z}.
\end{equation}
\begin{equation}
\begin{aligned}
    & \Pr(\{s_{g_Z}\}) = \frac{1}{2^{r_Z'} (2\cosh\beta)^{r_Z}} \mathcal{Z}(\{s_{g_Z}\})\\
    &= \frac{1}{2^{r_Z}}  \sum_{\sigma_Z \in \ker H_Z^T} (\tanh \beta)^{|\sigma_Z|} s_{\sigma_Z}
\end{aligned}
\end{equation}
We want to prove $\lim_{n\rightarrow\infty} I_c(\rho_L, \mathcal{N} \circ \mathcal{E}_s)/k < \log 2 $ given that $\ker H_Z^T$ is finite dimensional, $\dim \ker H_Z^T = r_Z - r_Z'= O(1)$.

\begin{equation}
    \begin{aligned}
        &\rho_0(s) =  \frac{(2\cosh\beta)^{r_Z}}{\mathcal{Z}(s)}\sum_{s_e|s_e+s\in \im H_Z}  (1-q)^{r_Z - s_e} q^{s_e} \ket{\sigma_{Z} = s+s_e}\bra{\sigma_{Z} = s+s_e} \otimes \ket{\sigma_{X}' = 0}\bra{\sigma_{X}' = 0} \otimes \Gamma\\
        &=  \frac{1}{\mathcal{Z}(s)}\sum_{\sigma_{Z}\in \im H_Z}  \exp\left(\sum_{e_2}\beta (-1)^{{s}_{e_2}+ \sigma_{Z,e_2}}\right) \ket{\sigma_{Z}}\bra{\sigma_{Z}}\otimes \ket{\sigma_{X}' = 0}\bra{\sigma_{X}' = 0} \otimes \Gamma\\
        &=  \frac{1}{\mathcal{Z}(s)}\sum_{s_e|s_e+s\in \im H_Z}  \exp\left(\sum_{e_2}\beta (-1)^{{s_e}_{,e_2}}\right) \ket{\sigma_{Z} = s+s_e}\bra{\sigma_{Z} = s+s_e}  \otimes \ket{\sigma_{X}' = 0}\bra{\sigma_{X}' = 0} \otimes \Gamma
    \end{aligned}
    \end{equation}
Note that if $\sigma_Z \notin \im H_Z$, then $\ket{\sigma_Z}\bra{\sigma_Z}  = 0$.
The post-noise state is 
\begin{equation}
    \begin{aligned}
        &\mathcal N(\rho_0(s)) =\sum_c \sum_{\sigma_{Z}\in \im H_Z}  \Pr(c) \Pr(\sigma_Z) X_c \left[\ket{\sigma_{Z}}\bra{\sigma_{Z}} \otimes \ket{\sigma_{X}' = 0}\bra{\sigma_{X}' = 0} \otimes \Gamma \right] X_c\\
    \end{aligned}
\end{equation}
where 
\begin{equation}
    \Pr(\sigma_Z) = \frac{1}{\mathcal{Z}(s)} \exp\left(\sum_{e_2}\beta (-1)^{{s}_{e_2}+ \sigma_{Z,e_2}}\right) = \frac{(2\cosh \beta)^{r_Z}}{\mathcal{Z}(s)} (1-q)^{r_Z -|s+\sigma_Z|} q^{|s+\sigma_Z|}.
\end{equation}
We purify the state as 
\begin{equation}
    \ket{\Psi_{RQAE}} = \frac{1}{\sqrt{2^k}} \sum_c \sum_{[h] \in H^1} \sum_{\sigma_{Z}\in \im H_Z}  \sqrt{\Pr(c) \Pr(\sigma_Z)} X_c \ket{\sigma_{Z}}\otimes \ket{\sigma_{X}' = 0} \otimes \ket{[h]} \otimes \ket{\sigma_Z}_A \otimes \ket{c}_E\otimes \ket{[h]}_R.
\end{equation}
Since
\begin{equation}
    -I_c +k \log 2 = - S(RAE) + S(AE) + S(R) = S(\rho_{RAE}||\rho_{AE} \otimes \rho_{R}) \geq \frac{1}{2\log 2} \norm{\rho_{RAE} - \rho_{AE} \otimes \rho_{R}}^2_1,
\end{equation}
define Schatten $\alpha$-norm as 
\begin{equation}
    \norm{X}_{\alpha} = \left(\tr \left[\left(\sqrt{X^\dagger X}\right)^\alpha\right]\right)^{\frac{1}{\alpha}},
\end{equation}
we have $\norm{X}_{\alpha} \leq \norm{X}_{1}$ for $\alpha \in (1,+\infty)$, thus
\begin{equation}
    S(\rho_{RAE}||\rho_{AE} \otimes \rho_{R})\geq \frac{1}{2\log 2} \norm{\rho_{RAE} - \rho_{AE} \otimes \rho_{R}}^2_2.
\end{equation}
Since every $X$ error chain $c$ is determined by a $Z$ syndrome $H_Z c$, a logical class and a $X$ stabilizer, we denote the logical class of c as $l(c) \in H^1$.
We have
\begin{equation}
    \begin{aligned}
        &\rho_{RAE} = \sum_{c_1,c_2 \in C^1} \sum_{\sigma_1,\sigma_2 \in \im H_Z} \sum_{[h_1],[h_2] \in H^1}\frac{1}{2^k} \delta(\sigma_1 + \sigma_2 + H_Z(c_1+c_2)= 0)\delta([h_1+h_2] + l(c_1+c_2) = 0)\\
        &\times \sqrt{\Pr(c_1) \Pr(c_2) \Pr(\sigma_1) \Pr(\sigma_2)} \ket{c_1}_E\bra{c_2}_E\otimes \ket{\sigma_1}_A\bra{\sigma_2}_A\otimes \ket{[h_1]}_R\bra{[h_2]}_R,
    \end{aligned}
\end{equation}
\begin{equation}
    \begin{aligned}
        &\rho_{AE} = \sum_{c_1,c_2 \in C^1} \sum_{\sigma_1,\sigma_2 \in \im H_Z} \frac{1}{2^k} \delta(\sigma_1 + \sigma_2 + H_Z(c_1+c_2)= 0)\delta( l(c_1+c_2) = 0)\\
        &\times \sqrt{\Pr(c_1) \Pr(c_2) \Pr(\sigma_1) \Pr(\sigma_2)} \ket{c_1}_E\bra{c_2}_E\otimes \ket{\sigma_1}_A\bra{\sigma_2}_A,
    \end{aligned}
\end{equation}
and $\rho_R = I_R/2^k$.
\begin{equation}
    \begin{aligned}
        &\rho_{RAE}- \rho_{AE} \otimes \rho_{R} = \frac{1}{2^k}\sum_{c_1,c_2 \in C^1} \sum_{\sigma_1,\sigma_2 \in \im H_Z}  \delta(\sigma_1 + \sigma_2 + H_Z(c_1+c_2)= 0)\sqrt{\Pr(c_1) \Pr(c_2) \Pr(\sigma_1) \Pr(\sigma_2)} \ket{c_1}_E\bra{c_2}_E\otimes \ket{\sigma_1}_A\bra{\sigma_2}_A\\
        &\otimes\sum_{[h_1],[h_2] \in H^1}\left\{ \delta([h_1+h_2] + l(c_1+c_2) = 0)- \delta([h_1]=[h_2])\delta( l(c_1+c_2) = 0)\right\}\ket{[h_1]}_R\bra{[h_2]}_R,
    \end{aligned}
\end{equation}

\begin{equation}
    \begin{aligned}
        &\norm{\rho_{RAE}- \rho_{AE} \otimes \rho_{R}}^2_2 = \tr [(\rho_{RAE}- \rho_{AE} \otimes \rho_{R})^2]\\  
        &=\frac{1}{2^{2k}} \sum_{c_1,c_2 \in C^1} \sum_{\sigma_1,\sigma_2 \in \im H_Z}  \delta(\sigma_1 + \sigma_2 + H_Z(c_1+c_2)= 0)\Pr(c_1) \Pr(c_2) \Pr(\sigma_1) \Pr(\sigma_2) \\
        &\times\sum_{[h_1],[h_2] \in H^1}\left\{ \delta([h_1+h_2] + l(c_1+c_2) = 0)- \delta([h_1]=[h_2])\delta( l(c_1+c_2) = 0)\right\}^2,\\
        &=\frac{1}{2^{2k}} \sum_{c_1,c_2 \in C^1} \sum_{\sigma_1,\sigma_2 \in \im H_Z}  \delta(\sigma_1 + \sigma_2 + H_Z(c_1+c_2)= 0)\Pr(c_1) \Pr(c_2) \Pr(\sigma_1) \Pr(\sigma_2) \\
        &\times\sum_{[h_1],[h_2] \in H^1}\left\{ \delta([h_1+h_2] + l(c_1+c_2) = 0)- \delta([h_1]=[h_2])\delta( l(c_1+c_2) = 0)\right\},\\
        &=\frac{1}{2^{k}} \sum_{c_1,c_2 \in C^1} \sum_{[h] \in H^1:[h]\neq 0}\sum_{\sigma_1,\sigma_2 \in \im H_Z}  \\
        &\times \delta(\sigma_1 + \sigma_2 + H_Z(c_1+c_2)= 0) \delta([h] = l(c_1+c_2))\Pr(c_1) \Pr(c_2) \Pr(\sigma_1) \Pr(\sigma_2) \\
        &=\frac{1}{2^{k}} \sum_{c_1,c_2 \in C^1} \sum_{[h] \in H^1:[h]\neq 0}\sum_{\sigma_1,\sigma_2 \in \im H_Z} \delta((\sigma_1)'=(\sigma_2)')  \delta([h] = l(c_1+c_2))\Pr(c_1) \Pr(c_2) \Pr(\sigma_1+H_Z c_1) \Pr(\sigma_2+H_Z c_2) \\
        &=\frac{1}{2^{k}} \sum_{c_1,c_2 \in C^1} \sum_{[h] \in H^1:[h]\neq 0}\sum_{\sigma \in \im H_Z}   \delta([h] = l(c_1+c_2))\Pr(c_1) \Pr(c_2) \Pr(\sigma+H_Z c_1) \Pr(\sigma +H_Z c_2) 
    \end{aligned}
\end{equation}
Since $\ker H_Z^T = (\im H_Z)^\perp$, we have
\begin{equation}
    \mathbb{I}_{\im H_Z} (x) = \frac{1}{|\ker H_Z^T|} \sum_{v \in \ker H_Z^T} (-1)^{\braket{v,x}}, 
\end{equation}

\section{Soundness guarantees single-shot preparation}
\subsection{Notations and preliminary lemmas}
For stochastic data errors $e\in C^1$ (or syndrome errors $s_e \in C^2$),  we denote the corresponding random variable using the capital letter $E$ ($S_e$). The probability of a specific configuration is denoted as 
\begin{equation}
    \Pr(E = e) = p(e).
\end{equation}
$\textrm{supp}(e)$ is the support of $e$, which is a set of data qubits ( $\textrm{supp}(s_e)$ for check bits). In the following parts, without leading to conflict, we use $e$ to represent both the error chain and its support set in order to simplify the notation. The probability that a set is contained in some error $E$ is denoted as 
\begin{equation}
    \Pr(e \subseteq E) = \sum_{e '\supseteq e} \Pr(E=e').
\end{equation}
\begin{definition}[Local stochastic error]
    We say a random error $E$ is $(p)$-local stochastic if there exists a parameter $p$ such that 
    \begin{equation}
        \Pr(e \subseteq E) \leq O\left(p^{|e|} \right).
    \end{equation}
\end{definition}

For an observed syndrome configuration $s \in C^2$, it might not stay in $\im H_Z$, so it is possible to be a nonphysical syndrome. So we need to project it back to the physical syndrome space $\im H_Z$. 
Notice that the physical syndrome space defined a classical linear code $\im H_Z \subseteq C^2$ in the full syndrome space, which is usually called 'meta check'.
The projection is performed by a minimum-weight decoder for syndrome errors. 
\begin{definition}[minimum weight decoder for syndrome error]
    Given the classical linear code $\im H_Z \subseteq C^2$, the minimum weight decoder $r: C^2 \rightarrow C^2$ is defined as follows: 
    \begin{equation}
        r(s) = \argmin_{s' \in V:s'+s\in \im H_Z}|s'|, \quad \forall s\in C^2.
    \end{equation}
\end{definition}
After decoding, we have the residual syndrome
\begin{equation}
    \sigma = \pi(s) = s + r (s) \in \im H_z.
\end{equation}
Under the minimum weight decoder, the observed syndrome $s$ must occupy more than half of the residual syndrome $\sigma$. Moreover, this fact holds for any physical sub-syndrome of $\sigma$.
\begin{lemma} \label{lemma:half-weight}
    Given $\sigma = \pi(s) = s+r(s)$, if $b \subseteq \sigma$ and $b \in \im H_Z$, then $|s \cap b| \geq |b|/2$.
\end{lemma}
\begin{proof}
    Because $r$ is the minimum weight decoder, $|r(s) \leq r(s)+b|$.
    So
    \begin{equation}
        |r(s)+b| = |r(s)| + |b| - 2|r(s) \cap b| \geq |r(s)|, 
    \end{equation}
    and thus $|r(s)\cap b| \leq |b|/2$.
    Denote the restriction of $\sigma$ on the support of $b$ as $\sigma|_b$. Since $\sigma|_b = \vec 1$, we have
    \begin{equation}
        s|_b = (\sigma + r(s))|_b = 1 + r(s)|_b, 
    \end{equation}
    and thus 
    \begin{equation}
        |\sigma \cap b| = |b| - |r(s) \cap b| \geq \frac{|b|}{2}.
    \end{equation}
\end{proof}


For a physical syndrome $\sigma$, we can decode the data error using the minimum weight decoder of (Pauli $X$) data error.
\begin{definition}[minimum weight decoder for data error]
    Given the party check matrix $H_Z$, the minimum weight decoder $c: \im H_Z \rightarrow C^1$ is defined as follows: 
    \begin{equation}
        c(\sigma) = \argmin_{c \in C^1: H_Z c = \sigma}|c|, \quad \forall \sigma \in \im H_Z.
    \end{equation}
\end{definition}
In this section, we will assume $H_Z$ is an LDPC code with check degree $d_c$ and variable degree $d_v$. Moreover, we require it to have linear soundness.
\begin{definition}[$(t,f)$-soundness]
    Let $t\in \mathbb{Z}$ and $f : \mathbb{Z} \rightarrow \mathbb{R}$ be a function with $f (0) = 0$. We say the parity check $H_Z$ is $(t,f)$ sound if for all physical syndromes $\sigma \in \im H_z$ with $\sigma \leq t $ it follows that $f(|\sigma|) \geq |c(\sigma)|$.
\end{definition}
\begin{definition}[Linear soundness]
    For a $(t,f)$-sound parity check, we say it is $(t,\rho)$-linear sound if there exists a constant $\rho > 0$ such that $f(x) = x/ \rho$.
\end{definition}
Notice that the definition is strictly weaker than soundness in Ref. \cite{quintavalleSingleShotErrorCorrection2021} since it bounds the weights of reduced errors corresponding to syndrome $\sigma$ in all logical classes.

Since $H_z$ is LDPC, we can define its bounded-degree adjacency graph as follows,
\begin{definition}[Adjacency graph]
Given a parity check matrix $H_Z$, we define the adjacency graph $G(H_Z)$ where we associate each of the $n$ qubits with a vertex and two qubits are connected by an edge if there is a check in $H_Z$ that acts nontrivially on both of them.    
\end{definition}
For $(d_v,d_c)$-LPDC $H_Z$, $G$ has a bounded vertex degree of at most $\Delta = (d_c-1)d_v$. In this paper, we always assume $\Delta > 1$.
To examine whether a certain random error variable $E$ is local stochastic,
for a set of vertices $e$ (or a data error configuration), we often need to consider the errors containing $e$, $e\subset E$. In the proof of the main theorem, we need to mediate this event by considering the connected components of $E$ that contain $e$. 
\begin{definition}[$e$-anchored]
Given two set of vertices $w$ and $e$, we say $w$ is $e$-anchored if $e \subseteq w$ and for the connected component of $w$, $w = w_1 \sqcup w_2 \cdots w_m$, we have $w_i \cap e \neq \emptyset$, $\forall w_i$.
\end{definition}
Note that here we use $\sqcup$ to denote the disconnected union on the adjacency graph. If $w$ is $e$-anchored, then each connected component of $w$ must contain some connected components of $e$. 
There is a counting lemma for $e$-anchored sets,
\begin{lemma}[Cluster counting lemma (lemma 2 of Ref. \cite{gottesmanFaultTolerantQuantumComputation2014})]
\label{lemma:counting}
Let $M_\Delta(\omega,e)$ be the total number of $e$-anchored sets that have weight $\omega$, then 
\begin{equation}
    M_\Delta(\omega,e) \leq \e^{|e|-1}(\Delta \e)^{\omega - |e|},
\end{equation}
where $\e$ is the base of the natural logarithm.
\end{lemma}
Given a set $w$, we often need to consider the case where $w$ is a union of some connected components of $E$. If $E = \bigsqcup_{i \in M} E_i$, then we denote $w \sqsubseteq E$ if there exists an index subset $M' \subseteq M$ such that $w = \bigsqcup_{i \in M'} E_i$. If here $E$ is given by the minimum weight decoder of data error, we have the following lemma
\begin{lemma}
\label{lemma:UCC}
    Given the parity check $H_Z$ and $w \in C^1$, $\sigma \in \im H_Z$. If on the adjacency graph $G(H_Z)$ we have $w \sqsubseteq c(\sigma)$, then 
    \begin{enumerate}
        \item the support of $H_Z w$ is contained in the support of $\sigma$, $H_Z w \subseteq \Sigma$;
        \item the minimum weight correction of $H_Z w$ is $w$, $c(H_Z w) = w$.
    \end{enumerate}
\end{lemma}
\begin{proof}
    Suppose that $c(\sigma) = w \sqcup f$. Since $w$ and $f$ are disconnected, they cannot share a syndrome check, otherwise they will be connected by this check on $G(H_Z)$. Then because $\sigma = H_z w + H_z f$ and $H_zw \cap H_z f = \emptyset$, we must have $H_Z w \subseteq \sigma$.

    Suppose that $c(H_Z w) \neq w$, there exist $w' \in C^1$ such that $H_z w =H_z w'$ and $|w'| < |w|$. So
    \begin{equation}
        H_Z (w'+f) = H_Z(w+f) = H_Z(c(\sigma)) = \sigma.
    \end{equation}
    However,
    \begin{equation}
        |w'+f| \leq |w'|+|f| < |w|+|f| =|c(\sigma)|,
    \end{equation}
    which is in contradiction with that $c(\sigma)$ is the minimum weight decoder.
    So $c(H_Z w) = w$.
\end{proof}

\subsection{state preparation process}
Suppose that the readout suffer from a stochastic error $S_e$, which has the distribution $\Pr(S_e=s_e)=q(s_e)$. Then starting from $+$ product state, the post-measurement is
    \begin{equation}
    \begin{aligned}
        &\rho_Q = \sum_s \Pr(s) \ket{s}_c\bra{s}_c \otimes  \sum_{s_e|s_e+s\in \im H_Z}  \tilde{q}(s_e) \ket{\sigma_{Z} = s+s_e}\bra{\sigma_{Z} = s+s_e}\\& \otimes \ket{\sigma_{X}' = 0}\bra{\sigma_{X}' = 0} \otimes \ket{[h_X]=0}\bra{[h_X]=0},\\
        &= \sum_s \Pr(s) \ket{s}_c\bra{s}_c \otimes  \sum_{\sigma \in \im H_Z}  \tilde{q}(\sigma+s) \ket{\sigma_{Z} = \sigma}\bra{\sigma_{Z} = \sigma} \otimes \ket{\sigma_{X}' = 0}\bra{\sigma_{X}' = 0} \otimes \ket{[h_X]=0}\bra{[h_X]=0},
    \end{aligned}
    \end{equation}
where the probability of observing a particular syndrome $s$ is
\begin{equation}
    \Pr(S=s) = \frac{1}{r_Z'}\sum_{\sigma \in \im H_z} q(\sigma+s) = \frac{1}{r_Z'} \Pr(S_e \in [s]_{\im H_Z})
\end{equation}
Here $[s]_{\im H_Z}$ denotes the equivalent class in the quotient space $C^2/\im H_z$. It means that the probability of observed syndrome depends only on its equivalent class. Also, since the projection of unphysical syndrome vanishes $\Pi_{\sigma_Z \notin \im H_Z}=0$, the post measurement state containes only superposition of physical syndrome, following from the distribution
\begin{equation}
\begin{aligned}
    &\tilde{q}(s_e) = \frac{q(s_e)}{\sum_{\sigma \in \im H_z} q(\sigma+s)} = \frac{q(s_e)}{\Pr(S_e \in [s]_{\im H_Z})}.
\end{aligned}
\end{equation}
Provided with the information of $s$, we apply the minimum weight decoder to correct syndrome error, and then apply accordingly the controlled Pauli $X$ correction. For each $s$ we apply $X_{c(\pi(s))}$ to $\rho(s)$, which yields the state
    \begin{equation}
    \begin{aligned}
        & \frac{1}{2^{r_z'}}\sum_s  \ket{s}_c\bra{s}_c \otimes  \sum_{s_e|s_e+s\in \im H_Z}  q(s_e) X_{c(\pi(s))} \ket{\sigma_{Z} = s+s_e}\bra{\sigma_{Z} = s+s_e}  \otimes \ket{\sigma_{X}' = 0}\bra{\sigma_{X}' = 0} \otimes \ket{[h_X]=0}\bra{[h_X]=0} X_{c(\pi(s))},\\
        &=\frac{1}{2^{r_z'}}\sum_s  \ket{s}_c\bra{s}_c \otimes  \sum_{s_e|s_e+s\in \im H_Z}  q(s_e) \ket{\sigma_{Z} = r(s)+s_e}\bra{\sigma_{Z} = r(s)+s_e}  \otimes \ket{\sigma_{X}' = 0}\bra{\sigma_{X}' = 0} \otimes \ket{[h_X]=0}\bra{[h_X]=0},\\
        &=\frac{1}{2^{r_z'}}\sum_s  \ket{s}_c\bra{s}_c \otimes  \sum_{\sigma \in \im H_Z}  q(\sigma+r(s)) \ket{\sigma_{Z} = \sigma}\bra{\sigma_{Z} = \sigma}  \otimes \ket{\sigma_{X}' = 0}\bra{\sigma_{X}' = 0} \otimes \ket{[h_X]=0}\bra{[h_X]=0}
    \end{aligned}
    \end{equation}
We then discard the classical register, such that the final state is 
\begin{equation}
        \rho_{\text{final}}= \sum_{\sigma \in \im H_Z}  \Pr(\Sigma = \sigma)\ket{\sigma_{Z} = \sigma}\bra{\sigma_{Z} = \sigma}  \otimes \ket{\sigma_{X}' = 0}\bra{\sigma_{X}' = 0} \otimes \ket{[h_X]=0}\bra{[h_X]=0}.
\end{equation}
where $\Sigma$ follows from the distribution
\begin{equation}
     \Pr(\Sigma = \sigma) = \frac{1}{2^{r_z'}}\sum_s q(\sigma+r(s)) .
\end{equation}
In the expression of $\rho_{\text{final}}$, the superposition of different syndromes $\sigma$ can be converted into the superposition of Pauli $c$ strings with has the corresponding syndrome $H_Z c = \sigma$. The choice could be arbitrary since $X$ stabilizers and $X$ logical operators act trivially on the logical $+$ state. In this section, we fix the choice as the minimum weight decoding for data error $c(\sigma)$. In that case, we have
\begin{equation}
    \rho_{\text{final}} = \mathcal{S}(\ket{\{+\}}_L \bra{\{+\}}_L ),
\end{equation}
where
\begin{equation}
    \mathcal{S}(\rho ) = \sum_{\sigma \in \im H_Z} \Pr(\Sigma = \sigma)  X_{c(\sigma)} \rho X_{c(\sigma)}.
\end{equation}

To analyze the single-shot preparation threshold, we need to study $\mathcal{S}$. However, a subtlety is that the set of Kraus operators $\{\sqrt{\Pr(\Sigma = \sigma)} X_{c(\sigma)} \}_{\sigma}$ exactly satisfies KL condition, 
\begin{equation}
    \Pi X_{c(\sigma_1)} X_{c(\sigma_2)} \Pi = \delta_{\sigma_1,\sigma_2} \Pi,
\end{equation}
since for $\sigma_1 \neq \sigma_2$, $c(\sigma_1)+c(\sigma_2)$ always have nontrivial syndrome, $H_Z(c(\sigma_1)+c(\sigma_2) ) = \sigma_1+\sigma_2 \neq 0$, thus yields $0$ under code space projection. This holds for arbitrary choices of the map $c: \im H_Z \rightarrow C^1$. Thus, the crucial criterion for state preparation is whether it remains correctable in the presence of additional stochastic Pauli $ X$-data noise. We need to determine whether the channel $\mathcal{N} \circ \mathcal{S}$ has a QEC threshold. It suffices to show $\mathcal{N} \circ \mathcal{S}$ is local stochastic, which follows the probability
\begin{equation}
    \Pr(F=\eta) = \sum_{\sigma \in \im H_Z } \Pr(\Sigma = \sigma) \Pr(E = \eta + c_{\sigma}).
\end{equation}

\subsection{local stochastic property of error distributions}
Given that $\Pr(S_e=s_e)=q(s_e)$ is the local stochastic syndrome error with probability parameter $q$, we consider the random variable 
\begin{equation}
    \Sigma = \pi(S_e) = S_e + r(S_e). 
\end{equation}
Its probability is evaluated as follows. Consider the quotient space $C^2/\im H_z$, denote its elements as $[s]_{\im H_Z}$. For $s,s' \in [s]_{\im H_Z}$, we have $r(s) = r(s')$.
So $[s]_{\im H_Z}$ is composed by $r(s) +\sigma$, $\sigma\in \im H_Z$. 
Conversely, for every $[s]_{\im H_Z} \in C^2/\im H_Z$, $s = r(s) +\sigma$ is the unique element in $[s]_{\im H_Z}$ such that $\pi(s) = \sigma$.
So the probability of $\pi(S_e) = \sigma$ is evaluated by collecting the probability of $S_e = r(s) +\sigma$ in each equivalent class.
\begin{equation}
    \Pr(\Sigma = \sigma)=\Pr(\pi(S_e) = \sigma) = \sum_{[s]\in C^2/\im H_Z} \Pr(S_e = \sigma + r(s)) = \frac{1}{2^{r_z'}} \sum_{s\in C^2}q(\sigma+r(s)).
\end{equation}
We then prove an important lemma that $c(\Sigma)$ is local stochastic.
\begin{lemma}
\label{lemma:local-stochastic-syndrome}
    Given an $(\kappa n,\rho)$-linear sound $(d_v,d_c)$-LDPC $H_Z$, let $\Sigma = \pi(S_e) = S_e + r(S_e)$ where $S_e$ is a $(q)$-local stochastic syndrome error, then $c(\Sigma)$ is $(\tilde{q})$-local stochastic with $\tilde{q} = \e (2\sqrt q)^{\min\{\rho, \kappa\} }$, provided that $\tilde{q} < 1/\Delta$.
\end{lemma}
\begin{proof}
    We need to evaluate the subset probability $\Pr(e \subseteq c(\Sigma))$. We then enumerate all the $e$-anchored sets $w$, such that $w$ is a union of some connected components of $c(\Sigma)$ (denoted as $w \sqsubseteq c(\Sigma)$). Therefore, for the event $\{e \subseteq c(\Sigma)\}$ we have
    \begin{equation}
        \{e \subseteq c(\Sigma)\} \subseteq \bigcup_{w: \text{ $e$-anchored}} \{w \sqsubseteq c(\Sigma)\}.
    \end{equation}
    Then the union lemma tells us
    \begin{equation}
        \Pr(e \subseteq c(\Sigma)) \leq \sum_{w: \text{ $e$-anchored}} \Pr(w \sqsubseteq c(\Sigma)).
    \end{equation}
    According to Lemma \ref{lemma:UCC}, if $w \sqsubseteq c(\Sigma)$ then $H_Z w \subseteq \Sigma$, so 
    \begin{equation}
        \Pr(e \subseteq c(\Sigma)) \leq \sum_{w: \text{ $e$-anchored}} \Pr(H_Z w \subseteq \Sigma).
    \end{equation}
    Since $\Sigma = \pi(S_e)$, for $H_Z w\subseteq \Sigma$, according to Lemma \ref{lemma:half-weight} we have $|S_e \cap H_Z w| \geq |H_Z w |/2$, therefore the event $\{H_Z w \subseteq \Sigma\}$ can be enumerated by the events that $S_e$ contains more than half of $H_Z w$,
    \begin{equation}
        \{H_Z w \subseteq \Sigma\} \subseteq \bigcup_{s \subseteq H_Z w: |s| \geq \lceil |H_z w |/2\rceil} \{s \subseteq S_e\}.
    \end{equation}
    So using the union lemma,
    \begin{equation}
    \begin{aligned}
        &        \Pr(H_Z w \subseteq \Sigma) \leq \sum_{s \subseteq H_Z w: |s| \geq \lceil |H_z w |/2\rceil} \Pr(s \subseteq S_e) \leq \sum_{s \subseteq H_Z w: |s| \geq \lceil |H_z w |/2\rceil} q^{|s|}\\
        & \leq \sum_{x = \lceil |H_z w |/2\rceil}^{|H_Z w|} \binom{|H_Zw|}{x} q^x \leq (2\sqrt{q})^{|H_Z w|}.
    \end{aligned}
    \end{equation}
    Since $H_z$ is $(t,\rho)$-sound with $t=\Theta(n)$, assume $t \geq \kappa n$,
    we have $|H_Z w| \geq \min \{\rho |w|, t \}$,
    \begin{equation}
        \Pr(H_Z w \subseteq \Sigma) \leq (2\sqrt{q})^{\min\{\rho|w|, \kappa n\}}.
    \end{equation}
    We can then evaluate
    \begin{equation}
    \begin{aligned}
        &\Pr(e \subseteq c(\Sigma)) \leq \sum_{w: \text{ $e$-anchored}} (2\sqrt{q})^{\min\{\rho|w|, \kappa n\}} = \sum_{x \geq |e|} M_\Delta (x,e) (2\sqrt{q})^{\min\{\rho x, \kappa n\}}\\
        & =\sum_{|e| \leq x \leq \kappa n/\rho} M_\Delta (x,e) (2\sqrt{q})^{\rho x} + (2\sqrt{q})^{\kappa n} \sum_{x > \kappa n/\rho} M_\Delta (x,e) \\
        & \leq \sum_{|e| \leq x \leq \kappa n/\rho} \e^{|e|-1}(\Delta \e)^{x - |e|} (2\sqrt{q})^{\rho x} + (2\sqrt{q})^{\kappa n}\sum_{ \kappa n/\rho <x \leq n} \e^{|e|-1}(\Delta \e)^{x - |e|}\\
        &=\e^{|e|-1}(2\sqrt q)^{\rho |e|}
        \frac{1-\big[\Delta \e(2\sqrt q)^\rho\big]^{\,\lfloor \kappa n/\rho\rfloor-|e|+1}}
        {1-\Delta \e(2\sqrt q)^\rho}+\e^{|e|-1}(2\sqrt q)^{\kappa n}
        (\Delta \e)^{\lfloor \kappa n/\rho\rfloor+1-|e|}
        \frac{(\Delta \e)^{\,n-\lfloor \kappa n/\rho\rfloor}-1}{\Delta \e-1}\\
        & \leq
        \frac{ (\e (2\sqrt q)^{\rho})^{|e|}}
        {\e(1-\Delta \e(2\sqrt q)^\rho)}+
        \frac{(\Delta \e) (\e (2\sqrt q)^\kappa)^{|e|}}{\e(\Delta \e-1)} \leq O\left( \left(\e (2\sqrt q)^{\min\{\rho, \kappa\} } \right)^{|e|}\right)
    \end{aligned}
    \end{equation}
    Where we assumed $\Delta \e(2\sqrt q)^\rho < 1$ and $\Delta \e(2\sqrt q)^\kappa < 1$.
\end{proof}


The next lemma states that the composition of two local stochastic errors is still local stochastic.
    \begin{lemma}
        Let $E_1$ be a $(p_1)$-local stochastic error and $E_2$ be a $(p_2)$-local stochastic error, then $E_1+E_2$ is a $(p_1 + p_2)$-local stochastic error.
    \end{lemma}
    \begin{proof}
        For any error configuration $e$ we have
        \begin{equation}
        \begin{aligned}
            &\Pr(e \subseteq E_1 + E_2) \leq \sum_{e' \subseteq e} \Pr(e' \subseteq E_1) \Pr(e\setminus e' \subseteq E_2 )\leq const \times \sum_{e' \subseteq e} p_1^{|e'|} p_2^{|e|-|e'|} \leq O\left((p_1 +p_2)^{|e|} \right). 
        \end{aligned}
        \end{equation}
    \end{proof}



\section{Statistical Mechanical Mapping}
In this section, we assume $c: \im H_Z \rightarrow C^1$ to be an arbitrary map satisfying $H_Zc(\sigma)=\sigma$. We assume that after the noise $\mathcal{N} \circ \mathcal{S}$ we apply a round of perfect error correction, that is, measure the syndrome of the residual error $F$ perfectly and apply the correction class with the largest probability. To find the optimal correction, we note that the decoding of the preparation round $r:C^2 \rightarrow C^2$ is done by a Pauli $X$ operator that commutes with the final correction, so we can postpone it and choose the correction according to both $s$ and $H_Z F$ in the final round. The post-measurement state would be
\begin{equation}
    \frac{1}{2^{r_z'}}\sum_{s\in C^2, \eta \in C^1}  \ket{s}_c\bra{s}_c \otimes \ket{H_Z \eta}_c\bra{H_Z \eta}_c \otimes  \sum_{\sigma \in \im H_Z}  q(\sigma+s) p(\eta + c(\sigma)) X_\eta \rho X_\eta.
\end{equation}
We denote the syndrome of $\eta$ as $\sigma_\eta = H_Z \eta$. By fixing the choice of  $c(\sigma_\eta)$, we can decompose $\eta = c(\sigma_\eta)+b+l$, $b \in B^1$, $[l] \in H^1$. Since both the choice of $c(\sigma)$ and $c(\sigma_\eta)$ are arbitrary, we can absorb the redefinition of $c(\sigma)$ into $c(\sigma_\eta)$. Consequently we have
\begin{equation}
    \sum_{s\in C^2, \sigma_\eta \in \im H_Z}  \ket{s}_c\bra{s}_c \otimes \ket{\sigma_\eta }_c\bra{\sigma_\eta }_c \otimes \sum_{[l]\in H^1} \Pr(s,\sigma_\eta,[l]) X_{c(\sigma_\eta)} X_l\rho X_{c(\sigma_\eta)} X_l,
\end{equation}
Where 
\begin{equation}
    \Pr(s,\sigma_\eta,[l]) = \frac{1}{2^{r_z'}}\sum_{\sigma\in \im H_Z, b\in B^1} q(\sigma+s)p(c(\sigma)+c(\sigma_\eta)+b+l).
\end{equation}
The maximum likelihood decoder (MLD) first fixes the syndrome part by applying $X_{c(\sigma_\eta)}$ and then applies the logical operator with the largest probability conditioned on $s$ and $\sigma_\eta$,
\begin{equation}
    l_{MLD}(s,\sigma_\eta) = \arg \max_{[l] \in H^1} \Pr([l]|s,\sigma_\eta).
\end{equation}
We define the disordered partition function as 
\begin{equation}
    \mathcal Z(s,\sigma_\eta)  =  \sum_{c\in C^1} q(H_Z c +s)p(c+\eta)  \propto \sum_{[l]} \Pr(s,\sigma_\eta,[l]),
\end{equation}
and view $c$ as the statistical mechanical d.o.f.. The disorder probability could be chosen as $\Pr(s,\sigma_\eta)$. To simplify the problem, according to the Nishimori physics, choosing a disorder probability proportional to the partition function is equivalent to choosing a disorder probability proportional to the Boltzmann weight. So we choose
\begin{equation}
    \Pr_{\text{dis}}(s,\eta) = q(s)p(\eta).
\end{equation}
The average free energy will be the same
\begin{equation}
\begin{aligned}
    &\mathcal F = - \sum_{s,\sigma_\eta} \Pr(s,\sigma_\eta)\log Z(s,\sigma_\eta)\\
    & \propto - \sum_{s,\eta} \sum_{c'\in C^1} q(H_Z c' +s)p(c'+\eta)\log \sum_{c\in C^1} q(H_Z c +s)p(c+\eta)\\
    &(s\rightarrow s+H_Z c',\eta\rightarrow \eta+c') \propto - \sum_{s,\eta} \sum_{c'\in C^1} q( s)p(\eta)\log \sum_{c\in C^1} q(H_Z (c+c') +s)p(c+c'+\eta)\\
    &(c\rightarrow c+c') \propto - \sum_{s,\eta}  q( s)p(\eta)\log \sum_{c\in C^1} q(H_Z c +s)p(c+\eta)\\
    &\propto - \sum_{s,\eta} \Pr_{\text{dis}}(s,\eta)\log Z(s,\sigma_\eta).
\end{aligned}
\end{equation}
We will consider single-qubit noise
\begin{equation}
    q(s) = q^{|s|}(1-q)^{r_Z-|s|}, \qquad p(\eta) = p^{|\eta|}(1-p)^{n-|\eta|}.
\end{equation}

\section{KW duality}
Let
\begin{equation}\label{eq:Zseta-csum}
Z_{s,\eta}
=\sum_{c\in C^1}
\left(\frac{q}{1-q}\right)^{|Hc+s|}
\left(\frac{p}{1-p}\right)^{|c+\eta|}.
\end{equation}
Introducing the Fourier representation of the binary delta function,
\begin{equation}\label{eq:binary-delta}
\delta(x)=2^{-r_Z}\sum_{\tau\in C^2}(-1)^{\langle \tau,x\rangle},
\end{equation}
we obtain
\begin{equation}\label{eq:Zseta-fourier}
Z_{s,\eta}
=2^{-r_Z}\sum_{\tau\in C^2}\sum_{\sigma\in C^2,c\in C^1}
\left(\frac{q}{1-q}\right)^{|\sigma+s|}
\left(\frac{p}{1-p}\right)^{|c+\eta|}
(-1)^{\langle \tau,\sigma+Hc\rangle}.
\end{equation}
After shifting variables \(u=\sigma+s\) and \(d=c+\eta\), and using
\begin{align}
\sum_{u\in C^2} a^{|u|}(-1)^{\langle \tau,u\rangle}
&=(1+a)^{m-|\tau|}(1-a)^{|\tau|},
\label{eq:hadamard-u}
\\
\sum_{d\in C^1} b^{|d|}(-1)^{\langle H^T\tau,d\rangle}
&=(1+b)^{n-|H^T\tau|}(1-b)^{|H^T\tau|},
\label{eq:hadamard-d}
\end{align}
with \(a=\frac{q}{1-q}\) and \(b=\frac{p}{1-p}\), one finds
\begin{equation}\label{eq:Zseta-dual-intermediate}
Z_{s,\eta}
=
2^{-r_Z}(1-q)^{-r_Z}(1-p)^{-n}
\sum_{\tau\in C^2}
(-1)^{\langle \tau,s\rangle+\langle H^T\tau,\eta\rangle}
(1-2q)^{|\tau|}(1-2p)^{|H^T\tau|}.
\end{equation}
Finally, using \(\langle H^T\tau,\eta\rangle=\langle \tau,H\eta\rangle\), this becomes
\begin{equation}\label{eq:Zseta-KW}
Z_{s,\eta}
=
2^{-r_Z}(1-q)^{-r_Z}(1-p)^{-n}
\sum_{\tau\in C^2}
(-1)^{\langle \tau,s+H\eta\rangle}
(1-2q)^{|\tau|}(1-2p)^{|H^T\tau|},
\end{equation}
which is the desired KW-dual form.

\section{Order Parameter}
To diagnose the MLD threshold, it is natural to decompose the partition function into logical sectors,
\begin{equation}
    \mathcal Z(s,\sigma_\eta)=\sum_{[l]\in H^1}\mathcal Z_{[l]}(s,\sigma_\eta),\qquad 
    \mathcal Z_{[l]}(s,\sigma_\eta):=\sum_{\sigma\in \im H_Z}\sum_{b\in B^1} q(\sigma+s)\,p\!\left(c(\sigma)+c(\sigma_\eta)+b+l\right),
\end{equation}
and to define the posterior distribution on \(H^1\cong \mathbb F_2^k\) by
\begin{equation}
    P_{s,\sigma_\eta}([l]):=\Pr([l]\mid s,\sigma_\eta)=\frac{\mathcal Z_{[l]}(s,\sigma_\eta)}{\mathcal Z(s,\sigma_\eta)}.
\end{equation}
After choosing a basis of logical \(X\)-classes and a dual basis of logical \(Z\)-cycles, every \(u\in H_1\cong \mathbb F_2^k\) defines a logical Wilson loop
\begin{equation}
    W_u([l]):=\chi_u([l])=(-1)^{u\cdot [l]},
\end{equation}
which depends only on the logical class \([l]\), not on the microscopic representative. Its thermal expectation value in the disorder sample \((s,\sigma_\eta)\) is
\begin{equation}
    m_u(s,\sigma_\eta):=\langle W_u\rangle_{s,\sigma_\eta}
    =\sum_{[l]\in H^1} \chi_u([l])\,P_{s,\sigma_\eta}([l])
    =\frac{1}{\mathcal Z(s,\sigma_\eta)}\sum_{[l]\in H^1}\chi_u([l])\,\mathcal Z_{[l]}(s,\sigma_\eta).
\end{equation}
Thus the family \(\{m_u\}_{u\in H_1}\) is precisely the Fourier transform of the posterior on logical sectors; in particular, the ordered phase relevant for decoding is the phase in which the Gibbs measure is polarized into one logical sector, equivalently \(m_u\to 1\) for all \(u\neq 0\) after a suitable choice of the reference sector, whereas in the disordered phase \(m_u\to 0\) for all \(u\neq 0\). A convenient gauge-invariant scalar order parameter is
\begin{equation}
    q_{\rm top}(s,\sigma_\eta):=\frac{1}{2^k-1}\sum_{u\neq 0} m_u(s,\sigma_\eta)^2.
\end{equation}

\begin{theorem}[Logical Wilson loops and MLD success probability]
Let \(k=\dim H^1\), and let
\begin{equation}
    P_{\rm succ}(s,\sigma_\eta):=\max_{[l]\in H^1} P_{s,\sigma_\eta}([l])
\end{equation}
be the success probability of the maximum-likelihood decoder for a fixed disorder sample \((s,\sigma_\eta)\). Then:
\begin{enumerate}
    \item The posterior is recovered exactly from the logical Wilson loop expectation values via
    \begin{equation}\label{eq:fourier-inversion-posterior}
        P_{s,\sigma_\eta}([l])=\frac{1}{2^k}\sum_{u\in H_1} \chi_u([l])\,m_u(s,\sigma_\eta).
    \end{equation}
    Consequently,
    \begin{equation}\label{eq:exact-success}
        P_{\rm succ}(s,\sigma_\eta)
        =\frac{1}{2^k}\max_{[l]\in H^1}\sum_{u\in H_1} \chi_u([l])\,m_u(s,\sigma_\eta).
    \end{equation}
    \item If \(q_{\rm top}(s,\sigma_\eta)=\frac{1}{2^k-1}\sum_{u\neq 0}m_u(s,\sigma_\eta)^2\), then
    \begin{equation}\label{eq:success-bounds}
        \frac{1+(2^k-1)q_{\rm top}(s,\sigma_\eta)}{2^k}
        \;\le\;
        P_{\rm succ}(s,\sigma_\eta)
        \;\le\;
        \sqrt{\frac{1+(2^k-1)q_{\rm top}(s,\sigma_\eta)}{2^k}}.
    \end{equation}
\end{enumerate}
In particular, \(q_{\rm top}\to 0\) implies \(P_{\rm succ}\to 2^{-k}\), while \(q_{\rm top}\to 1\) implies \(P_{\rm succ}\to 1\).
\end{theorem}

\begin{proof}
Since \(H^1\cong \mathbb F_2^k\) is a finite abelian group, its characters satisfy the orthogonality relation
\begin{equation}
    \frac{1}{2^k}\sum_{u\in H_1} \chi_u([l])\chi_u([l'])=\delta_{[l],[l']}.
\end{equation}
Using the definition \(m_u=\sum_{[l']} \chi_u([l'])P_{s,\sigma_\eta}([l'])\), we obtain
\begin{equation}
    \frac{1}{2^k}\sum_{u} \chi_u([l])\,m_u
    =\frac{1}{2^k}\sum_{u}\sum_{[l']} \chi_u([l])\chi_u([l'])P_{s,\sigma_\eta}([l'])
    =\sum_{[l']} \delta_{[l],[l']} P_{s,\sigma_\eta}([l'])
    =P_{s,\sigma_\eta}([l]),
\end{equation}
which proves \eqref{eq:fourier-inversion-posterior}. Taking the maximum over \([l]\) immediately gives \eqref{eq:exact-success}. Next, Parseval's identity on the finite group \(H^1\) yields
\begin{equation}
    \sum_{[l]\in H^1} P_{s,\sigma_\eta}([l])^2
    =\frac{1}{2^k}\sum_{u\in H_1} m_u(s,\sigma_\eta)^2
    =\frac{1+(2^k-1)q_{\rm top}(s,\sigma_\eta)}{2^k},
\end{equation}
since \(m_0=1\). For any probability distribution \(P\) on a finite set,
\begin{equation}
    \sum_x P(x)^2 \le \max_x P(x)\le \sqrt{\sum_x P(x)^2},
\end{equation}
because the left inequality follows from \(P(x)^2\le P(x)\max_y P(y)\) and summing over \(x\), while the right inequality is immediate from \(\max_x P(x)^2\le \sum_x P(x)^2\). Applying this to \(P_{s,\sigma_\eta}\) gives \eqref{eq:success-bounds}. The final claims follow by taking the limits \(q_{\rm top}\to 0\) and \(q_{\rm top}\to 1\).
\end{proof}

\begin{proposition}
Let $H^1\simeq \mathbb F_2^k$ be the logical sector group, and set
\[
N:=|H^1|=2^k.
\]
For each disorder sample $x=(s,\eta)$, let
\begin{equation}\label{eq:posterior-sector}
P_x([l]) := \Pr([l]\mid s,\sigma_\eta), \qquad [l]\in H^1,
\end{equation}
be the posterior distribution on logical sectors, and define the samplewise MLD success probability by
\begin{equation}\label{eq:psucc-sample}
P_{\mathrm{succ}}(x):=\max_{[l]\in H^1} P_x([l]).
\end{equation}
For each character $u\in \widehat{H^1}\cong \mathbb F_2^k$, define the logical Wilson expectation
\begin{equation}\label{eq:logical-wilson}
m_u(x):=\sum_{[l]\in H^1} (-1)^{u\cdot [l]} P_x([l]),
\end{equation}
and the topological Edwards--Anderson parameter
\begin{equation}\label{eq:qtop-def}
q_{\mathrm{top}}(x):=\frac{1}{N-1}\sum_{u\neq 0} m_u(x)^2.
\end{equation}
Then, for every disorder sample $x$,
\begin{equation}\label{eq:sample-two-sided}
\frac{1+(N-1)q_{\mathrm{top}}(x)}{N}
\;\le\;
P_{\mathrm{succ}}(x)
\;\le\;
\sqrt{\frac{1+(N-1)q_{\mathrm{top}}(x)}{N}}.
\end{equation}
Consequently, after averaging over disorder,
\begin{equation}\label{eq:disorder-two-sided}
\frac{1+(N-1)[q_{\mathrm{top}}]_{\mathrm{dis}}}{N}
\;\le\;
[P_{\mathrm{succ}}]_{\mathrm{dis}}
\;\le\;
\sqrt{\frac{1+(N-1)[q_{\mathrm{top}}]_{\mathrm{dis}}}{N}},
\end{equation}
where
\[
[P_{\mathrm{succ}}]_{\mathrm{dis}}:=\sum_x \Pr_{\mathrm{dis}}(x)\,P_{\mathrm{succ}}(x),
\qquad
[q_{\mathrm{top}}]_{\mathrm{dis}}:=\sum_x \Pr_{\mathrm{dis}}(x)\,q_{\mathrm{top}}(x).
\]
Equivalently, since $N=2^k$,
\begin{equation}\label{eq:disorder-two-sided-2k}
\frac{1+(2^k-1)[q_{\mathrm{top}}]_{\mathrm{dis}}}{2^k}
\;\le\;
[P_{\mathrm{succ}}]_{\mathrm{dis}}
\;\le\;
\sqrt{\frac{1+(2^k-1)[q_{\mathrm{top}}]_{\mathrm{dis}}}{2^k}}.
\end{equation}
\end{proposition}

\begin{proof}
Fix a disorder sample $x=(s,\eta)$ and abbreviate
\[
P([l]) := P_x([l]), \qquad m_u:=m_u(x), \qquad q_{\mathrm{top}}:=q_{\mathrm{top}}(x).
\]
Since $P$ is a probability distribution on the finite abelian group $H^1$, Parseval's identity on $H^1$ gives
\begin{equation}\label{eq:parseval-proof}
\sum_{[l]\in H^1} P([l])^2
=
\frac{1}{N}\sum_{u\in \widehat{H^1}} m_u^2.
\end{equation}
Moreover, by definition of the trivial character,
\[
m_0=\sum_{[l]\in H^1}P([l])=1.
\]
Therefore, using \eqref{eq:qtop-def},
\begin{equation}\label{eq:l2-via-qtop}
\sum_{[l]\in H^1} P([l])^2
=
\frac{1}{N}\left(1+\sum_{u\neq 0}m_u^2\right)
=
\frac{1+(N-1)q_{\mathrm{top}}}{N}.
\end{equation}

Now let
\[
P_{\max}:=\max_{[l]\in H^1}P([l])=P_{\mathrm{succ}}(x).
\]
For any probability distribution on $N$ points one has
\begin{equation}\label{eq:elementary-bounds}
\sum_{[l]\in H^1}P([l])^2 \le P_{\max} \le \sqrt{\sum_{[l]\in H^1}P([l])^2}.
\end{equation}
Indeed, the left inequality follows from
\[
\sum_{[l]}P([l])^2 \le P_{\max}\sum_{[l]}P([l])=P_{\max},
\]
while the right inequality follows from
\[
P_{\max}^2 \le \sum_{[l]}P([l])^2.
\]
Substituting \eqref{eq:l2-via-qtop} into \eqref{eq:elementary-bounds} yields the samplewise bound \eqref{eq:sample-two-sided}.

Finally, average \eqref{eq:sample-two-sided} over the disorder distribution. The lower bound averages directly:
\[
[P_{\mathrm{succ}}]_{\mathrm{dis}}
\ge
\frac{1+(N-1)[q_{\mathrm{top}}]_{\mathrm{dis}}}{N}.
\]
For the upper bound, by concavity of $\sqrt{\cdot}$ and Jensen's inequality,
\[
[P_{\mathrm{succ}}]_{\mathrm{dis}}
\le
\sum_x \Pr_{\mathrm{dis}}(x)\sqrt{\frac{1+(N-1)q_{\mathrm{top}}(x)}{N}}
\le
\sqrt{\frac{1+(N-1)[q_{\mathrm{top}}]_{\mathrm{dis}}}{N}}.
\]
This proves \eqref{eq:disorder-two-sided}, and \eqref{eq:disorder-two-sided-2k} is just the specialization $N=2^k$.
\end{proof}
\subsection{A local logical order parameter}

The natural order parameter for the MLD threshold is the polarization of the Gibbs measure among logical sectors. 
To realize this polarization by an observable that factorizes over sites, it is convenient to choose a \emph{linear} section
\begin{equation}\label{eq:linear-section}
    r:\im H_Z \to C^1,
    \qquad
    H_Z r(\sigma)=\sigma,
\end{equation}
which always exists after choosing a linear complement of \(\ker H_Z\subset C^1\). 
For a fixed disorder sample \((s,\eta)\), with \(\sigma_\eta:=H_Z\eta\), define the syndrome-corrected total error
\begin{equation}\label{eq:ell-r-def}
    \ell_r(c;\eta):=c+\eta+r\!\bigl(H_Z(c+\eta)\bigr)
    =c+\eta+r(H_Zc+\sigma_\eta).
\end{equation}
By construction,
\begin{equation}\label{eq:ell-r-cocycle}
    H_Z\ell_r(c;\eta)=H_Z(c+\eta)+H_Zr(H_Z(c+\eta))=0,
\end{equation}
hence \(\ell_r(c;\eta)\in Z^1\), so its homology class \([\ell_r(c;\eta)]\in H^1=Z^1/B^1\) is well defined. 
Now let \(u\in H_1\cong \mathbb F_2^k\) be a nontrivial logical character, and choose a representative \(\bar z_u\in C_1\) satisfying
\begin{equation}\label{eq:dual-logical-rep}
    \langle \bar z_u,b\rangle =0 \quad \forall\, b\in B^1,
    \qquad
    \langle \bar z_u,l\rangle =u\cdot [l] \quad \forall\, l\in Z^1 .
\end{equation}
Equivalently, \(\bar z_u\) represents a dual logical \(Z\)-operator. 
We then define the local logical order parameter
\begin{equation}\label{eq:local-order-parameter}
    O_u(c;\eta):=(-1)^{\langle \bar z_u,\ell_r(c;\eta)\rangle}
    =(-1)^{u\cdot [\ell_r(c;\eta)]}.
\end{equation}
Since \(r\) is linear, \eqref{eq:ell-r-def} becomes
\begin{equation}\label{eq:ell-r-linear}
    \ell_r(c;\eta)=(I+rH_Z)c+(\eta+r\sigma_\eta),
\end{equation}
and therefore
\begin{equation}\label{eq:local-order-factorized}
    O_u(c;\eta)
    =(-1)^{\langle \bar z_u,\eta+r\sigma_\eta\rangle}
     (-1)^{\langle (I+rH_Z)^T\bar z_u,\,c\rangle}
    =\xi_u(\eta)\prod_{i=1}^n (-1)^{\Lambda_{u,i}c_i},
\end{equation}
where
\begin{equation}\label{eq:Lambda-def}
    \Lambda_u:=(I+rH_Z)^T\bar z_u\in C_1,
    \qquad
    \xi_u(\eta):=(-1)^{\langle \bar z_u,\eta+r\sigma_\eta\rangle}.
\end{equation}
Thus \(O_u(c;\eta)\) indeed factorizes over sites, up to the disorder-dependent Mattis gauge factor \(\xi_u(\eta)\). Note that this factor is canceled in $q_{\text{top}}$. So we can equivalently consider the order parameter
\begin{equation}
    O_u(c)=(-1)^{\langle (I+rH_Z)^T\bar z_u,\,c\rangle}
    =\prod_{i=1}^n (-1)^{\Lambda_{u,i}c_i}.
\end{equation}

The corresponding thermal average is
\begin{equation}\label{eq:m-u-local}
    m_u(s,\eta)
    :=\frac{1}{\mathcal Z(s,\sigma_\eta)}
    \sum_{c\in C^1}
    O_u(c;\eta)\,q(H_Zc+s)\,p(c+\eta).
\end{equation}
Since \(O_u(c;\eta)\) depends only on the logical class \([\ell_r(c;\eta)]\), if we denote by
\begin{equation}\label{eq:sector-posterior-local}
    P_{s,\sigma_\eta}([l])
    :=\Pr([l]\mid s,\sigma_\eta)
\end{equation}
the posterior distribution of the logical class under the Gibbs measure, then grouping the sum in \eqref{eq:m-u-local} by logical sectors gives
\begin{equation}\label{eq:m-u-fourier}
    m_u(s,\eta)
    =\sum_{[l]\in H^1} (-1)^{u\cdot [l]} P_{s,\sigma_\eta}([l]).
\end{equation}
Hence \(\{m_u\}_{u\in H_1}\) is precisely the Fourier transform of the posterior distribution on \(H^1\).


Equation \eqref{eq:local-order-factorized} shows that \(O_u(c;\eta)\) is a genuine site-factorized observable, of Mattis type, whose thermal expectation detects whether the Gibbs measure is polarized in a definite logical sector. Therefore the MLD threshold can be studied as an ordering transition of the local observables \(O_u(c;\eta)\), or equivalently of the scalar quantity \(q_{\mathrm{top}}(s,\eta)\).

\subsection{Changing the section without changing the MLD success probability}

The definition of the logical-sector weights involves a choice of section
\begin{equation}\label{eq:section-c-def}
    c:\im H_Z\to C^1,
    \qquad
    H_Z c(\sigma)=\sigma .
\end{equation}
A different choice of representatives does not necessarily change the decoding problem. 
The correct invariance is under adding boundaries pointwise.

\begin{proposition}\label{prop:boundary-shift-invariance}
Let \(c:\im H_Z\to C^1\) be a section as in \eqref{eq:section-c-def}, and let
\begin{equation}\label{eq:r-boundary-shift}
    r:\im H_Z\to C^1,
    \qquad
    r(\sigma)=c(\sigma)+\beta(\sigma),
    \qquad
    \beta(\sigma)\in B^1
\end{equation}
for every \(\sigma\in \im H_Z\). Then \(r\) is also a section,
\begin{equation}\label{eq:r-is-section}
    H_Z r(\sigma)=\sigma,
\end{equation}
and the logical-sector weights defined using \(r\),
\begin{equation}\label{eq:Zl-r-def}
    \mathcal Z^{(r)}_{[l]}(s,\sigma_\eta)
    :=\sum_{\sigma\in \im H_Z}\sum_{b\in B^1}
    q(\sigma+s)\,
    p\!\left(r(\sigma)+r(\sigma_\eta)+b+l\right),
\end{equation}
coincide exactly with those defined using \(c\),
\begin{equation}\label{eq:Zl-c-def}
    \mathcal Z^{(c)}_{[l]}(s,\sigma_\eta)
    :=\sum_{\sigma\in \im H_Z}\sum_{b\in B^1}
    q(\sigma+s)\,
    p\!\left(c(\sigma)+c(\sigma_\eta)+b+l\right).
\end{equation}
Consequently,
\begin{equation}\label{eq:posterior-invariance}
    \Pr_r([l]\mid s,\sigma_\eta)=\Pr_c([l]\mid s,\sigma_\eta)
\end{equation}
for all \([l]\in H^1\), and therefore the MLD success probability is unchanged:
\begin{equation}\label{eq:success-invariance}
    P_{\mathrm{succ}}^{(r)}(s,\sigma_\eta)
    =
    P_{\mathrm{succ}}^{(c)}(s,\sigma_\eta).
\end{equation}
\end{proposition}

\begin{proof}
Since \(B^1\subset \ker H_Z\), we have
\begin{equation}\label{eq:section-proof-1}
    H_Z r(\sigma)=H_Z c(\sigma)+H_Z\beta(\sigma)=\sigma,
\end{equation}
so \(r\) is a section. Next, substituting \eqref{eq:r-boundary-shift} into \eqref{eq:Zl-r-def} gives
\begin{equation}\label{eq:section-proof-2}
\begin{aligned}
    \mathcal Z^{(r)}_{[l]}(s,\sigma_\eta)
    &=\sum_{\sigma\in \im H_Z}\sum_{b\in B^1}
    q(\sigma+s)\,
    p\!\left(c(\sigma)+\beta(\sigma)+c(\sigma_\eta)+\beta(\sigma_\eta)+b+l\right) \\
    &=\sum_{\sigma\in \im H_Z}\sum_{b\in B^1}
    q(\sigma+s)\,
    p\!\left(c(\sigma)+c(\sigma_\eta)+\bigl(b+\beta(\sigma)+\beta(\sigma_\eta)\bigr)+l\right).
\end{aligned}
\end{equation}
For each fixed \(\sigma\), the map
\begin{equation}\label{eq:b-translation}
    b\longmapsto b+\beta(\sigma)+\beta(\sigma_\eta)
\end{equation}
is a bijection of \(B^1\). Hence the sum over \(b\in B^1\) is unchanged, and
\begin{equation}\label{eq:section-proof-3}
    \mathcal Z^{(r)}_{[l]}(s,\sigma_\eta)
    =
    \mathcal Z^{(c)}_{[l]}(s,\sigma_\eta).
\end{equation}
Dividing by the total partition function \(\mathcal Z(s,\sigma_\eta)=\sum_{[l]}\mathcal Z_{[l]}(s,\sigma_\eta)\) proves \eqref{eq:posterior-invariance}, and \eqref{eq:success-invariance} follows by taking the maximum over \([l]\in H^1\).
\end{proof}

Proposition~\ref{prop:boundary-shift-invariance} shows that one may freely replace \(c(\sigma)\) by any representative in the same affine set
\begin{equation}\label{eq:affine-coset}
    c(\sigma)+B^1
\end{equation}
without changing the posterior on logical sectors or the MLD success probability. 
A convenient canonical construction is obtained by fixing a direct-sum decomposition
\begin{equation}\label{eq:direct-sum-U}
    C^1=B^1\oplus U
\end{equation}
and letting
\begin{equation}\label{eq:pi-U-def}
    \pi_U:C^1\to U
\end{equation}
be the linear projection along \(B^1\). Then one may define
\begin{equation}\label{eq:r-canonical}
    r(\sigma):=\pi_U\bigl(c(\sigma)\bigr).
\end{equation}
Since \(c(\sigma)-r(\sigma)\in B^1\), this choice satisfies the hypothesis of Proposition~\ref{prop:boundary-shift-invariance}, and therefore leaves the MLD success probability unchanged. 
In practice, one can choose \(U\) by Gaussian elimination on a basis of \(B^1\), so that the projection \(\pi_U\) is efficiently computable and provides a deterministic gauge-fixing of the section \(c\).

\subsection{A linear section \(r:\im H_Z\to C^1\) suitable for classical simulation.}
Let \(H_Z\in \mathbb F_2^{m\times n}\) be the parity-check matrix, and let \(\rho=\rank H_Z\). 
A convenient, exact construction of a linear section
\begin{equation}\label{eq:r-section-def}
    r:\im H_Z\to C^1=\mathbb F_2^n,
    \qquad
    H_Z r(\sigma)=\sigma \quad \forall\, \sigma\in \im H_Z,
\end{equation}
is obtained by Gaussian elimination over \(\mathbb F_2\). 
More precisely, one may choose an invertible matrix \(E\in GL(m,\mathbb F_2)\) and a permutation matrix \(\Pi\in GL(n,\mathbb F_2)\) such that
\begin{equation}\label{eq:HZ-systematic}
    E H_Z \Pi
    =
    \begin{pmatrix}
        I_\rho & A\\
        0 & 0
    \end{pmatrix},
\end{equation}
for some matrix \(A\in \mathbb F_2^{\rho\times (n-\rho)}\). 
For any \(\sigma\in \im H_Z\), define
\begin{equation}\label{eq:y-def}
    y:=E\sigma.
\end{equation}
Since \(\sigma\in \im H_Z\), the last \(m-\rho\) components of \(y\) vanish automatically, so \(y\) has the form
\begin{equation}\label{eq:y-split}
    y=
    \binom{\hat \sigma}{0},
    \qquad
    \hat \sigma\in \mathbb F_2^\rho.
\end{equation}
We then define the section \(r\) by setting all non-pivot variables to zero:
\begin{equation}\label{eq:r-def-systematic}
    r(\sigma):=
    \Pi
    \binom{\hat \sigma}{0}
    \in \mathbb F_2^n.
\end{equation}
This map is linear, since it is obtained by composition of linear maps, and it satisfies
\begin{equation}\label{eq:r-right-inverse}
    H_Z r(\sigma)=\sigma
    \qquad \forall\, \sigma\in \im H_Z.
\end{equation}
Indeed,
\begin{equation}\label{eq:r-right-inverse-proof}
    E H_Z r(\sigma)
    =
    E H_Z \Pi \binom{\hat \sigma}{0}
    =
    \begin{pmatrix}
        I_\rho & A\\
        0 & 0
    \end{pmatrix}
    \binom{\hat \sigma}{0}
    =
    \binom{\hat \sigma}{0}
    =
    E\sigma,
\end{equation}
and since \(E\) is invertible, \eqref{eq:r-right-inverse} follows.

This construction is particularly suitable for classical simulation, because the expensive preprocessing step---Gaussian elimination producing \(E\) and \(\Pi\)---is performed only once. 
After that, evaluating \(r(\sigma)\) for many syndromes requires only:
\begin{enumerate}
    \item computing \(y=E\sigma\),
    \item reading off the first \(\rho\) entries \(\hat \sigma\),
    \item padding with \(n-\rho\) zeros,
    \item and applying the inverse column permutation \(\Pi\).
\end{enumerate}
Thus the section can be applied efficiently in Monte Carlo sampling or exact summation routines. 
In practice, one should store the row-reduction producing \(E\) as a sequence of xor operations, rather than as a dense matrix, so that applying \(E\) to a syndrome is computationally cheap. 
If \(H_Z\) is sparse, a sparsity-aware pivot choice is preferable to reduce fill-in. 
We stress, however, that while this \(r\) is exact and efficient for classical computation, it is generally not geometrically local.
\bibliography{ref}
\end{document}